\clearpage
\thispagestyle{empty}
\pdfbookmark[0]{Berichtigung der Schneiden einer Waage. Ausgabe und Quellen}{balance-edition}
\begin{center}\Large\textbf{Berichtigung der Schneiden einer Waage}\end{center}

Der vollständige Bericht über Carl Friedrich Gauß’ Verfahren zur Justierung der Schneiden einer Hebelwaage ist eine KI-Übertragung des deutschen Textes in den \textit{Werken}, Band V (Göttingen, 1867), S.~511--513, in modernes Deutsch. Der Satz wurde mit KI erstellt.

Der Bericht erschien in den \textit{Göttingischen gelehrten Anzeigen}, 41.~Stück, 13. März 1837, S.~401--405. Der Nachdruck in den \textit{Astronomischen Nachrichten}, Bd.~14, Nr.~327 (1837), Sp.~241--244, nennt Gauß als Verfasser. Der Text spricht von Gauß in der dritten Person und berichtet über seinen Vortrag in der Königlichen Gesellschaft der Wissenschaften.

\medskip\noindent
\href{https://gdz.sub.uni-goettingen.de/id/PPN236006339}{Deutsche Werkausgabe (1867)}\par
\noindent\href{https://api.digitale-sammlungen.de/iiif/image/v2/bsb10538714_00407/full/1150,/0/default.jpg}{Originalausgabe der Zeitschrift, S.~401}\par
\noindent\href{https://mctoon.net/wp-content/uploads/2019/12/astronomische-nachrichten-313-336.pdf\#page=141}{Nachdruck in den Astronomischen Nachrichten (1837)}\par
\noindent\href{https://doi.org/10.1007/978-3-642-49319-5_26}{Deutscher Nachdruck (1877)}
\clearpage
